\documentclass[10pt,a4paper]{article}

\usepackage[T1]{fontenc}
\usepackage[utf8]{inputenc}
\usepackage[slovak]{babel}
\usepackage{lmodern}
\usepackage[margin=2cm]{geometry}
\usepackage{microtype}
\usepackage{array}
\usepackage{tabularx}
\usepackage{booktabs}
\usepackage{enumitem}
\usepackage{xcolor}
\usepackage{fancyhdr}
\usepackage{lastpage}
\usepackage[hidelinks]{hyperref}

\hypersetup{
  pdftitle={Black Box -- formuláre VV-A až VV-E},
  pdfauthor={Kyryl Tkachenko, Tymofii Trembovetskyi, Alex Vyšný}
}

\setlength{\parindent}{0pt}
\setlength{\parskip}{0.4em}
\renewcommand{\arraystretch}{1.3}

% údaje, ktoré treba doplniť ručne
\newcommand{\dopln}[1]{\textcolor{red}{[#1]}}

\pagestyle{fancy}
\fancyhf{}
\lhead{\small VV 2026 \quad Aplikovaný výskum}
\rhead{\small Akronym: BlackBox}
\cfoot{\small Strana \thepage{} z~\pageref{LastPage}}

% hlavička časti formulára
\newcommand{\cast}[2]{%
  \par\medskip
  \noindent\colorbox{black!10}{\parbox{\dimexpr\linewidth-2\fboxsep}{%
    \textbf{#1}\hfill\textbf{#2}}}\par\smallskip}

% riadok formulára: kód položky | názov | obsah
\newenvironment{polozky}
  {\noindent\tabularx{\linewidth}{|>{\centering\arraybackslash}p{7mm}|>{\raggedright\arraybackslash}p{45mm}|X|}\hline}
  {\endtabularx\par}
\newcommand{\pol}[3]{#1 & #2 & #3 \\ \hline}

\begin{document}

\begin{center}
  {\Large\bfseries Žiadosť o~poskytnutie prostriedkov na aplikovaný výskum}\\[0.4em]
  {\large Formuláre VV-A až VV-E}\\[0.4em]
  Black Box: lokálna aplikácia na včasné varovanie študentov pred študijným preťažením\\[0.3em]
  {\small Kyryl Tkachenko, Tymofii Trembovetskyi, Alex Vyšný \quad\textbar\quad
  Metódy inžinierskej práce, FIIT STU \quad\textbar\quad 30. október 2026}
\end{center}

%=====================================================================
\cast{VV-A1}{Základné informácie o~projekte}

\begin{polozky}
\pol{01}{Evidenčné číslo projektu}{APVV-26-xxxx (pridelí agentúra)}
\pol{02}{Dátum podania}{30.~10.~2026}
\pol{03}{Názov projektu}{Black Box: lokálna aplikácia na včasné varovanie študentov pred študijným preťažením}
\pol{03}{Názov projektu v~angličtine}{Black Box: A~Local Application for Early Warning of Students Against Study Overload}
\pol{04}{Akronym projektu}{BlackBox}
\pol{05}{Odbor vedy a~techniky}{Technické vedy -- informatika (softvérové inžinierstvo, interakcia človeka s~počítačom)}
\pol{06}{Charakter výskumu}{aplikovaný výskum}
\pol{07}{Začiatok riešenia projektu}{01.~01.~2027}
\pol{08}{Koniec riešenia projektu}{31.~12.~2028}
\pol{10}{Žiadateľská organizácia}{Slovenská technická univerzita v~Bratislave, Fakulta informatiky a~informačných technológií}
\pol{11}{Požadované finančné prostriedky z~APVV (v~EUR)}{63\,274}
\pol{12}{Spolufinancovanie projektu (v~EUR)}{0}
\pol{13}{Celkové náklady na projekt (v~EUR)}{63\,274}
\end{polozky}

\textbf{VV-A1-09 \quad Anotácia} (max. 2\,000 znakov)

Vysokoškolskí študenti dostávajú termíny zadaní, testov a~skúšok z~viacerých
zdrojov -- z~informačného systému školy, z~e-mailov aj zo skupinových chatov.
Chýba im jeden prehľad, ktorý by včas ukázal, že sa v~niektorom týždni
nahromadí priveľa povinností. Bežné kalendáre termíny iba zobrazujú
a~aplikácie zamerané na duševnú pohodu zasa nepoznajú študijné povinnosti.
Dôsledkom je preťaženie, stres a~v~dlhodobom horizonte riziko vyhorenia.

Cieľom projektu je navrhnúť a~overiť desktopovú aplikáciu Black Box, ktorá
funguje ako „čierna skrinka“ študenta. Aplikácia zhromaždí termíny na jednom
mieste, a~to ručným zadaním alebo importom kalendára vo formáte iCalendar.
Ku každému termínu používateľ odhadne, koľko hodín mu zaberie. Z~týchto
odhadov, kalibrovaných podľa skutočnej náročnosti minulých termínov,
a~z~blízkosti termínov aplikácia vypočíta záťaž jednotlivých týždňov
a~7 až 14 dní vopred upozorní na týždeň, ktorý bude pravdepodobne preťažený.
Raz denne môže používateľ zaznamenať, ako sa cíti, na stupnici od 1 do 5.
Po náročnom období aplikácia zobrazí časovú os termínov, vypočítanej záťaže
a~zaznamenaných pocitov, z~ktorej je zrejmé, čo preťaženie spôsobilo.

Projekt sa zameriava najmä na: (1) návrh modelu výpočtu týždennej záťaže
a~prahu pre varovanie, (2) implementáciu aplikácie v~nástroji Tauri
s~jadrom v~jazyku Rust a~s~lokálnou databázou SQLite, (3) overenie so
skupinou najmenej 30 študentov počas semestra, či varovania zodpovedajú
skutočne náročným týždňom.

Aplikácia je navrhnutá podľa princípu local-first: funguje bez internetu
a~bez servera a~všetky údaje vrátane citlivých údajov o~pohode zostávajú
v~zariadení používateľa v~súlade s~nariadením GDPR. Výsledkom projektu bude
funkčný prototyp s~otvoreným zdrojovým kódom a~vyhodnotenie presnosti
varovaní.

%=====================================================================
\cast{VV-A2}{Základné informácie o~riešiteľských organizáciách}

\textbf{Žiadateľ}

\begin{polozky}
\pol{01}{Názov organizácie}{Slovenská technická univerzita v~Bratislave}
\pol{02}{Skrátený názov}{STU}
\pol{03}{Adresa organizácie}{Vazovova 5, 812~43 Bratislava; pracovisko: Fakulta informatiky a~informačných technológií STU, Ilkovičova 2, 842~16 Bratislava}
\pol{04}{IČO}{00397687}
\pol{05}{Príslušnosť k~rezortu}{Ministerstvo školstva, výskumu, vývoja a~mládeže SR}
\pol{06}{Forma hospodárenia}{verejná vysoká škola}
\pol{07}{Kontaktná osoba}{Tymofii Trembovetskyi, \texttt{xtrembovetskyi@stuba.sk}, tel.~\dopln{doplniť}}
\pol{08}{Štatutárny zástupca I}{prof. Dr. Ing. Oliver Moravčík, rektor}
\pol{09}{Štatutárny zástupca II}{--}
\end{polozky}

\textbf{Spoluriešiteľská organizácia}

Projekt nemá spoluriešiteľskú organizáciu.

%=====================================================================
\cast{VV-A3}{Zoznam riešiteľov}

\textbf{01 \quad Zoznam zamestnancov priamo sa podieľajúcich na riešení projektu}

\noindent
\begin{tabularx}{\linewidth}{|>{\raggedright\arraybackslash}p{40mm}|c|X|c|c|c|>{\raggedright\arraybackslash}p{22mm}|}
\hline
\textbf{Meno a~priezvisko} & \textbf{Tituly} & \textbf{Pracovné zaradenie} & \textbf{Dátum nar.} & \textbf{IČO org.} & \textbf{Hodín} & \textbf{Hodín v~rokoch} \\ \hline
Tymofii Trembovetskyi (zodpovedný riešiteľ) & -- & študent, riešiteľ & \dopln{doplniť} & 00397687 & 800 & 2027: 400; 2028: 400 \\ \hline
Kyryl Tkachenko & -- & študent, riešiteľ & \dopln{doplniť} & 00397687 & 800 & 2027: 400; 2028: 400 \\ \hline
Alex Vyšný & -- & študent, riešiteľ & \dopln{doplniť} & 00397687 & 800 & 2027: 400; 2028: 400 \\ \hline
\end{tabularx}

\begin{polozky}
\pol{02}{Ostatní zamestnanci}{celkový počet ostatných osôb: 0; súhrnná kapacita: 0 hodín}
\pol{03}{Spolu}{celkový počet zamestnancov: 3; súhrnná kapacita: 2\,400 hodín}
\end{polozky}

%=====================================================================
\cast{VV-A4}{Základné informácie o~zodpovednom riešiteľovi}

\begin{polozky}
\pol{01}{Meno a~priezvisko}{Tymofii Trembovetskyi}
\pol{02}{Pohlavie}{\dopln{doplniť}}
\pol{03}{Telefón}{\dopln{doplniť}}
\pol{04}{E-mail}{\texttt{xtrembovetskyi@stuba.sk}}
\pol{05}{Zodpovedný riešiteľ je mladým vedeckým pracovníkom}{áno}
\pol{06}{Vedecké databázy (typ, ID výskumníka)}{zatiaľ bez záznamu (ORCID bude založený na začiatku riešenia projektu)}
\pol{07}{Projekty zodpovedného riešiteľa za posledných 5~rokov}{Počet: 0. Zodpovedný riešiteľ je študentom a~doteraz neriešil projekt financovaný z~grantovej schémy.}
\pol{08}{Expertízy, konzultácie a~ostatné výsledky s~priamym využitím v~praxi za posledných 5~rokov}{Počet: 0.}
\pol{09}{Aplikačné výstupy -- chránené}{Počet: 0.}
\end{polozky}

\textbf{Čestné vyhlásenie zodpovedného riešiteľa}

Ja, dolu podpísaný Tymofii Trembovetskyi, čestne vyhlasujem, že údaje
uvedené v~častiach „Základné informácie o~zodpovednom riešiteľovi“
a~v~štruktúrovanom profesijnom životopise sú pravdivé.

Ja, dolu podpísaný Tymofii Trembovetskyi, čestne vyhlasujem, že v~rámci
riešenia projektu budú zo strany žiadateľa rešpektované všetky zásady
etických princípov vyplývajúce z~európskej a~národnej legislatívy.

\medskip
\begin{tabularx}{\linewidth}{XXX}
Miesto: Bratislava & Dátum: 30.~10.~2026 & \rule{45mm}{0.4pt}\newline Podpis zodpovedného riešiteľa
\end{tabularx}

%=====================================================================
\cast{VV-B}{Ciele, zámery a~výstupy projektu}

\begin{polozky}
\pol{01}{Kľúčové slová}{študijná záťaž; včasné varovanie; vyhorenie študentov; osobná informatika; local-first; ochrana súkromia}
\end{polozky}

\textbf{02 \quad Ciele projektu}

Hlavným cieľom projektu je navrhnúť, implementovať a~overiť lokálnu
aplikáciu, ktorá študentovi vopred ukáže preťažené týždne semestra
a~po ich uplynutí pomôže pochopiť ich príčiny. Hlavný cieľ napĺňajú
štyri čiastkové ciele s~merateľnými ukazovateľmi:

\begin{enumerate}[label=C\arabic*,leftmargin=*,nosep]
  \item \textbf{Model záťaže.} Navrhnúť model výpočtu týždennej záťaže
    z~odhadov hodín a~dátumov termínov s~kalibráciou odhadov a~pravidlom
    pre varovanie. \emph{Ukazovateľ:} presnosť varovaní aspoň 70\,\%
    a~úplnosť aspoň 60\,\% pri predstihu najmenej 7 dní; model prekoná
    východiskový model (počet termínov v~týždni) v~ukazovateli F1.
  \item \textbf{Kalibrácia odhadov.} Overiť, že kalibrácia podľa
    skutočnej náročnosti minulých termínov znižuje chybu odhadu.
    \emph{Ukazovateľ:} pokles strednej absolútnej chyby odhadu hodín
    po kalibrácii oproti nekalibrovaným odhadom.
  \item \textbf{Prototyp aplikácie.} Implementovať desktopovú aplikáciu
    pre Windows, Linux a~macOS, ktorá funguje bez pripojenia
    k~internetu. \emph{Ukazovatele:} správne načítanie aspoň 95\,\%
    udalostí z~exportu informačného systému vo formáte iCalendar;
    nulová odchádzajúca sieťová komunikácia s~osobnými údajmi.
  \item \textbf{Overenie v~praxi.} Overiť aplikáciu s~najmenej 30 študentmi
    počas 8 týždňov semestra. \emph{Ukazovatele:} priemerné skóre
    použiteľnosti SUS aspoň 71; denný záznam vyplnený aspoň v~50\,\%
    dní.
\end{enumerate}

\begin{polozky}
\pol{03}{Zaradenie projektu do TRL stupnice}{TRL 2 $\rightarrow$ TRL 6}
\pol{03}{Zdôvodnenie zaradenia projektu do TRL stupnice}{Na začiatku projektu je sformulovaný koncept aplikácie a~jej využitia (TRL~2). Projekt prejde overením modelu na simulovaných rozvrhoch (TRL~3--4), vývojom prototypu (TRL~5) až po overenie funkčného prototypu so študentmi počas reálneho semestra, teda v~relevantnom prostredí (TRL~6).}
\pol{03}{Zaradenie výstupov projektu do TRL stupnice}{Model záťaže a~pravidlo varovania: TRL~4 (overené na simulovaných a~historických údajoch). Aplikácia Black Box: TRL~6 (prototyp overený v~relevantnom prostredí -- u~študentov počas semestra). Metodika overovania presnosti varovaní: TRL~6.}
\pol{04}{Odberateľ výsledkov je žiadateľ}{áno -- FIIT STU môže aplikáciu ponúknuť svojim študentom, napríklad v~rámci adaptácie študentov prvého ročníka, a~využiť zistenia o~rozložení záťaže pri plánovaní termínov v~predmetoch.}
\pol{04}{Odberateľ výsledkov je iný odberateľ}{nie (aplikácia bude dostupná s~otvoreným zdrojovým kódom aj pre študentov iných vysokých škôl)}
\pol{04}{Iný odberateľ (názov organizácie)}{--}
\end{polozky}

\textbf{05 \quad Prehľad plánovaných verejne dostupných výstupov a~prínosov podľa etáp}

\noindent
\begin{tabularx}{\linewidth}{|>{\raggedright\arraybackslash}p{38mm}|X|}
\hline
\textbf{Etapa} & \textbf{Výstupy a~prínosy} \\ \hline
E1 Analýza a~model záťaže (01/2027--06/2027) & Technická správa s~opisom modelu záťaže, kalibrácie a~pravidla varovania; simulátor modelu s~otvoreným kódom. \\ \hline
E2 Vývoj prototypu (04/2027--01/2028) & Prototyp Black Box 0.1 (import iCalendar, ručné zadanie, výpočet záťaže) a~verzia 1.0 (varovania, denný záznam, časová os, export a~vymazanie údajov); zdrojový kód pod otvorenou licenciou. \\ \hline
E3 Pilotné overenie (02/2028--07/2028) & Protokol overenia, anonymizovaný súhrnný dataset ukazovateľov, správa o~presnosti varovaní a~použiteľnosti. \\ \hline
E4 Vyhodnotenie a~šírenie (06/2028--12/2028) & Upravená verzia 2.0 podľa výsledkov overenia; príspevok na vedeckej konferencii; workshop pre študentov; záverečná správa. \\ \hline
\end{tabularx}

%=====================================================================
\cast{VV-C}{Rozpočet projektu v~EUR}

Žiadateľ: Slovenská technická univerzita v~Bratislave

\noindent
\begin{tabularx}{\linewidth}{|c|X|r|r|r|}
\hline
 & \textbf{Položka} & \textbf{2027} & \textbf{2028} & \textbf{Celkovo} \\ \hline
01 & Bežné priame náklady & 26\,014 & 26\,714 & 52\,728 \\ \hline
02 & \quad Mzdové náklady a~ostatné osobné náklady & 16\,800 & 16\,800 & 33\,600 \\ \hline
03 & \quad Zdravotné a~sociálne poistenie & 5\,914 & 5\,914 & 11\,828 \\ \hline
04 & \quad Cestovné náklady & 300 & 1\,200 & 1\,500 \\ \hline
05 & \quad Materiál & 2\,400 & 200 & 2\,600 \\ \hline
06 & \quad Služby & 400 & 2\,000 & 2\,400 \\ \hline
07 & \quad Energie, vodné, stočné, komunikácie & 0 & 0 & 0 \\ \hline
08 & \quad Náklady na popularizáciu & 200 & 600 & 800 \\ \hline
09 & Bežné nepriame náklady (20\,\% z~priamych) & 5\,203 & 5\,343 & 10\,546 \\ \hline
10 & Bežné náklady spolu & 31\,217 & 32\,057 & 63\,274 \\ \hline
   & \textbf{Celkové náklady z~APVV} & \textbf{31\,217} & \textbf{32\,057} & \textbf{63\,274} \\ \hline
   & Spolufinancovanie (štátne, zahraničné, súkromné) & 0 & 0 & 0 \\ \hline
   & \textbf{Celkové náklady} & \textbf{31\,217} & \textbf{32\,057} & \textbf{63\,274} \\ \hline
\end{tabularx}

\textbf{Rozpis predpokladaných nákladov uplatňovaných z~APVV}

\begin{description}[leftmargin=0pt,itemsep=2pt]
  \item[02 Mzdové náklady.] Traja riešitelia, každý 400 hodín ročne pri
    hodinovej sadzbe 14~EUR: $3 \times 400 \times 14 = 16\,800$~EUR ročne.
  \item[03 Poistenie.] Odvody zamestnávateľa vo výške 35,2\,\% z~mzdových
    nákladov.
  \item[04 Cestovné náklady.] 2027: účasť na odbornom podujatí v~SR
    (300~EUR). 2028: aktívna účasť na zahraničnej konferencii
    s~prezentáciou výsledkov (1\,200~EUR).
  \item[05 Materiál.] 2027: dva testovacie notebooky (macOS a~Windows) na
    zostavovanie a~testovanie aplikácie na všetkých podporovaných
    platformách (2\,400~EUR). 2028: drobný materiál (200~EUR).
  \item[06 Služby.] Každoročne certifikáty na podpisovanie kódu a~vývojárske
    účty pre distribúciu aplikácie (400~EUR). 2028 navyše: odmeny pre
    účastníkov pilotného overenia vo forme poukážok, 40 osôb po 30~EUR
    (1\,200~EUR), a~konferenčný poplatok (400~EUR).
  \item[07 Energie.] Hradené z~nepriamych nákladov pracoviska.
  \item[08 Popularizácia.] 2027: webová stránka projektu (200~EUR).
    2028: workshop pre študentov a~propagačné materiály (600~EUR).
  \item[09 Nepriame náklady.] 20\,\% z~bežných priamych nákladov
    na prevádzku pracoviska a~administratívu.
\end{description}

%=====================================================================
\cast{VV-D}{Harmonogram projektu}

\noindent
\begin{tabularx}{\linewidth}{|c|X|c|c|}
\hline
 & \textbf{Názov etapy} & \textbf{Začiatok etapy} & \textbf{Koniec etapy} \\ \hline
E1 & Analýza požiadaviek, návrh a~simulačné overenie modelu záťaže & 01/2027 & 06/2027 \\ \hline
E2 & Vývoj prototypu aplikácie Black Box (verzie 0.1 a~1.0) & 04/2027 & 01/2028 \\ \hline
E3 & Pilotné overenie so študentmi počas letného semestra & 02/2028 & 07/2028 \\ \hline
E4 & Vyhodnotenie, úprava aplikácie (verzia 2.0) a~šírenie výsledkov & 06/2028 & 12/2028 \\ \hline
\end{tabularx}

\textbf{Zoznam prístrojov v~projekte}

Projekt nepožaduje nákup prístrojov ani zariadení, ktoré by bolo potrebné
evidovať ako dlhodobý majetok.

%=====================================================================
\cast{VV-E}{Čestné vyhlásenie štatutárneho zástupcu žiadateľskej organizácie}

Ja, dolu podpísaný prof. Dr. Ing. Oliver Moravčík, štatutárny zástupca
žiadateľa, záväzne vyhlasujem, že:

\begin{itemize}[nosep,leftmargin=*]
  \item organizácia má platné osvedčenie o~spôsobilosti vykonávať výskum
    a~vývoj v~zmysle § 26a zákona č.~172/2005 Z.\,z. v~znení neskorších
    predpisov a~je evidovaná v~zozname spôsobilých osôb vykonávať výskum
    a~vývoj;
  \item všetky informácie obsiahnuté v~dokumentácii návrhu sú pravdivé;
  \item predložený návrh projektu nebol a~nie je financovaný z~iných zdrojov
    (národných alebo zahraničných);
  \item organizácia neposkytne žiadnu nepriamu štátnu pomoc v~zmysle
    špecifických podmienok výzvy;
  \item organizácia nie je daňovým dlžníkom, dlžníkom poistného na
    sociálnom poistení, dlžníkom príspevkov na starobné dôchodkové
    sporenie ani dlžníkom poistného na zdravotnom poistení;
  \item organizácia nie je v~likvidácii a~neporušila zákaz nelegálnej práce
    a~nelegálneho zamestnávania;
  \item organizácia má vysporiadané finančné vzťahy so štátnym rozpočtom
    a~bude súčinná pri poskytovaní dokladov podľa zákona č.~523/2004 Z.\,z.;
  \item som si vedomý povinnosti poskytnúť informácie o~výskume a~vývoji
    financovanom zo štátneho rozpočtu na účely zverejnenia a~povinnosti
    popularizovať výsledky riešenia projektu;
  \item som si vedomý, že APVV zverejní podľa § 19d ods.~4 zákona
    č.~172/2005 Z.\,z. rozhodnutie o~žiadosti v~rozsahu: číslo žiadosti,
    názov projektu, žiadateľ;
  \item predložený návrh projektu po svojej obsahovej stránke nebude
    financovaný z~národných zdrojov iných grantových schém.
\end{itemize}

V~prípade vyzvania zo strany APVV sa zaväzujem nahradiť toto čestné
vyhlásenie aktuálnymi potvrdeniami príslušných úradov.

\medskip
\begin{tabularx}{\linewidth}{XX}
Názov žiadateľskej organizácie: & Slovenská technická univerzita v~Bratislave \\
V~zastúpení: & prof. Dr. Ing. Oliver Moravčík, rektor \\
Dátum: & 30.~10.~2026 \\[1.5em]
 & \rule{50mm}{0.4pt}\newline Podpis štatutárneho zástupcu, pečiatka \\
\end{tabularx}

Čestné vyhlásenie štatutárneho zástupcu spoluriešiteľskej organizácie sa
neprikladá, pretože projekt nemá spoluriešiteľskú organizáciu.

\end{document}
